% Generated by code/integrate.py; edit extraction rules there.
\section{Policy Evaluation, Approximation, and Complete Action Coverage}\label{app:r15evaluation}
% Reviewed source: ECTA.tex, line 138, commit f5021cefa71492babfcbfa580e0c984f59a9de26
The distinction changes the population problem. With one action, $f(v)=-v$, $T=1$, and terminal utility one, the correct solution is $v^*(t)=e^{t-1}$. The perturbation $v_\epsilon=v^*-\epsilon(1-t)$ preserves the terminal condition. The joint objective under joint differentiation satisfies
\begin{equation}\label{eq:counterexample}
 \Delta(L_A+L_C)=-\frac\epsilon2+\frac{7\epsilon^2}{3},
 \qquad \Delta L_C=\frac{7\epsilon^2}{3}.
\end{equation}
Thus the block-specific critic objective has the correct root as its minimizer along this perturbation, whereas the joint objective does not. The same problem is visible in a stationary equation: $-H+\omega H^2$ prefers $H=1/(2\omega)$, not $H=0$. No choice of finite positive weight resolves this identification error.

The actor objective is a local policy-improvement surrogate. Under differentiability of the policy value, the derivative of lifetime recursive utility from $(t,x)$ instead involves the adjoint weight
\begin{equation}\label{eq:gradient}
 D_\phi V^{a_\phi}(t,x)=\E_{t,x}\int_t^\tau
 e^{\int_t^s f_y(v,S_v,a_v,V^{a_\phi})dv}
 D_a\HH^{a_\phi}[V^{a_\phi}](s,S_s)D_\phi a_\phi(s,S_s)\,ds,
\end{equation}
subject to the differentiability and integrability conditions given in the supplement. Arbitrary domain sampling does not reproduce this occupation and recursive-utility weighting. NBO uses \eqref{eq:actorloss} for approximate pointwise improvement and verifies the remaining action gap; it does not identify that loss with an exact lifetime policy gradient.

% Reviewed source: ECTA.tex, line 171, commit f5021cefa71492babfcbfa580e0c984f59a9de26
\subsubsection{Conditioning and stopping}
A positive, parameter-independent state weighting can precondition the critic equation without changing its zero set. If that weight depends on learned quantities, its derivative must be included in the stated critic objective. The homothetic laboratory divides the Hamiltonian by its positive coefficient and differentiates that entire normalized critic residual. In the actor step the same coefficient is detached. Error conversion to the original economic equation includes the scale factor.

A zero residual target does not make learning dynamics stationary. For the scalar loss $\ell(z)=a^2z^2/2$, gradient descent is stable only when $0<\eta a^2<2$. Likewise, full differentiation of $\delta^2/2$ gives $\delta\nabla\delta$, while a semi-gradient TD update differentiates only the current-value term. The two computations must not be compared using inconsistent signs or discount conventions. If $f$ already includes $-\rho v$, a time discretization must not add a second independent discount to the same continuation equation.

Constant-step Adam is a numerical choice in the calibration laboratory, not an application of a two-timescale theorem. A genuine two-timescale analysis would require nonsummable and square-summable step sizes, an actor-to-critic ratio tending to zero, controlled noise, bounded iterates, and suitable attracting limiting dynamics \citep{borkar2024}. A stationary critic gradient need not imply a zero residual, and a stationary multi-agent gradient need not imply Nash optimality. The results below therefore concern verified approximation errors rather than an assumed universal optimizer convergence property.

% Reviewed source: supp.tex, line 58, commit f5021cefa71492babfcbfa580e0c984f59a9de26
\subsection{Policy Evaluation and the Recursive Policy Gradient}
Let $a_\phi$ be differentiable in a scalar parameter direction, and suppose $V^{a_\phi}$ is sufficiently differentiable for the following calculation. The terminal payoff is independent of the policy parameter. Assume differentiation under the expectation is justified, that the differentiated state and utility equations are integrable, and that the boundary conditions of the differentiated problem are zero. Denote the directional derivative of the policy value by $w$.

Differentiating the fixed-policy equation gives
\[
 0=w_t+\LL^{a_\phi}w+f_y(t,x,a_\phi,V^{a_\phi})w
  +D_a\HH^{a_\phi}[V^{a_\phi}]\,D_\phi a_\phi.
\]
The operator derivative includes the dependence of both drift and covariance on the action. The derivative of the generator acting on the changing value is contained in $\LL^{a_\phi}w$, while the action derivative of the generator acting on the fixed value is contained in $D_a\HH$. No Hessian-action channel is discarded.

Applying the linear Feynman--Kac formula to the differentiated equation gives
\[
 w(t,x)=\E_{t,x}^{a_\phi}\int_t^\tau
 \exp\left(\int_t^s f_y(v,S_v,a_v,V^{a_\phi})dv\right)
 D_a\HH^{a_\phi}[V^{a_\phi}](s,S_s)D_\phi a_\phi(s,S_s)ds.
\]
This formula concerns the performance derivative from the stated initial state. A domain-uniform Hamiltonian objective uses a different weighting. Such an objective is useful for approximate pointwise improvement, but it becomes a lifetime policy gradient only with the appropriate occupation and adjoint weights. Trajectories sampled by the current actor are detached in the implemented domain objective. Differentiating through their sampling law would define a different gradient estimator and must be analyzed separately.

The critic's residual derivative is the derivative of the entire fixed-policy equation. In particular, its level is determined by the recursive evaluation equation together with boundary data. A separate negative-Hamiltonian gradient in the critic is not an innocuous level anchor: in the scalar stationary example it moves the minimizer away from the Bellman root.

\subsection{Representation, Boundary Conditions, and Error Measurement}
\subsubsection{Smooth approximation and what it establishes}
A derivative-approximation result concerns a smooth target on a compact domain and a family of networks whose capacity can increase. It does not imply that a fixed network attains zero residual, that optimization finds its best approximation, or that a nonsmooth target has classical derivatives everywhere. In a smooth problem, increasing the representation class can reduce approximation error in the value and derivatives. In a nonsmooth problem, NBO instead evaluates a consistent positive-weight Bellman equation, or solves an explicitly regularized equation before taking a justified limit.

The hard-terminal architecture is $v(t,x)=G(x)+(T-t)N(t,x)$. It exactly enforces $v(T,x)=G(x)$ independently of training. A lateral Dirichlet face needs an additional compatible lifting and vanishing factor. For a reflected state, the Neumann condition is a different constraint. A loss on only interior points leaves all of these boundary questions unanswered. The regression test evaluates the terminal architecture at randomly chosen states and checks exact equality, not merely a small mean terminal error.

The recursive Epstein--Zin domain is $bV>0$. An exponential output ensures positivity, while an interval transformation can keep $bV$ between declared positive barriers. If $bG>0$, one terminal-compatible positive representation is
\[
 bv(t,x)=\exp\left[\log(bG(x))+(T-t)N(t,x)\right].
\]
Its terminal value is exact. It does not by itself impose a lateral exit payoff or a two-sided barrier. These additional conditions must be built into the architecture or measured explicitly. Clipping an invalid utility argument inside a fractional power changes the equation and is not an acceptable hidden numerical repair.

\subsubsection{From residual samples to a uniform bound}
Let $r$ be $K$-Lipschitz on a domain in $\R^d$. If every point is within distance $h_V$ of a validation point, the triangle inequality gives $\|r\|_\infty\leq\max_{x\in\mathcal V}|r(x)|+Kh_V$. If $K$ is estimated from samples rather than bounded, the resulting expression is itself a diagnostic, not a certified uniform bound.

An integral bound needs more assumptions. Suppose the sampling density is at least $p_0>0$, and each ball of radius $r_0$ or less centered in the domain intersects the domain in volume at least $c_Dr^d$. Let $M=\|r\|_\infty$ and assume $M/(2K)\leq r_0$. Near a point approaching the maximum, Lipschitz continuity gives $|r|\geq M/2$ on a ball of radius $M/(2K)$. Hence
\[
 \E r^2\geq\frac{p_0c_D}{2^{d+2}K^d}M^{d+2},\qquad
 M\leq\left(\frac{2^{d+2}K^d\E r^2}{p_0c_D}\right)^{1/(d+2)}.
\]
For larger $M$, use the capped radius $r_0$ and the corresponding two-case inequality. A concentration argument is additionally needed to replace the population mean by an empirical one. This calculation explains why a dimension-free uniform certificate cannot be read directly from an average training loss.

The action gap is treated similarly. If an action net has covering radius $h_A$ and the Hamiltonian is $K_A$-Lipschitz in the action, the unobserved supremum exceeds the net maximum by at most $K_Ah_A$. Thus the net gap plus this term is an upper bound. A collection of successful deviations without a covering or duality argument is only a lower bound on exploitability.

% Reviewed source: ECTA.tex, line 240, commit f5021cefa71492babfcbfa580e0c984f59a9de26
\subsubsection{Why an integrated differential residual does not select the value}
Consider the exit problem $F(v')=1-|v'|=0$ on $(-1,1)$ with zero boundary payoff. Its value is $v^*(x)=|x|-1$. The smooth functions
\[
 w_\epsilon(x)=\sqrt{1+\epsilon^2}-\sqrt{x^2+\epsilon^2}
\]
satisfy the same boundary values and have vanishing integrated squared residual, but converge uniformly to $1-|x|$. At zero the test function one touches the wrong limit from above and gives $F(0)=1>0$, violating the viscosity subsolution inequality. The defect is not removed by exact Hamiltonian maximization.

On a grid with spacing $h$, the positive-weight exit backup is
\begin{equation}\label{eq:exit}
 (Tv)_i=-h+\max(v_{i-1},v_{i+1}),\qquad v_0=v_N=0.
\end{equation}
It is an undiscounted shortest-exit problem, not an instance of the stationary contraction theorem. Backward propagation of the boundary values gives $v_i=|x_i|-1$ exactly. At the peak of the wrong candidate, the \emph{scaled} residual $|w_i-(Tw)_i|/h$ is close to two, rather than close to zero. The laboratory therefore distinguishes correct selection from small integrated differential error.

An explicit positive viscosity term is another possible regularization, provided comparison, boundary convergence, and approximation errors for the regularized equations are controlled before taking the viscosity to zero. Network width or spectral bias is not itself a positive viscosity parameter. Moreover, approximation of a $C^2$ function and its derivatives does not establish approximation of nonexistent classical derivatives at a kink.

\subsubsection{Coverage and independent maximization}\label{sec:coverage}
A verification distribution is declared independently of the actor's occupation measure. It includes economically relevant counterfactual states and boundary regions, not only states visited by a current policy. If a residual $r$ has a known Lipschitz bound $K$ on a compact domain and a validation set has fill distance $d_\mathrm{fill}$, then
\begin{equation}\label{eq:cover}
 \|r\|_\infty\leq\max_{x\in\mathcal V}|r(x)|+K d_\mathrm{fill}.
\end{equation}
The same bound can be applied to an action gap if its modulus is controlled. An empirical mean alone does not give this guarantee. If the sampling density is bounded below and the domain has a uniform interior-volume condition, a Lipschitz residual also admits an $L^2$-to-sup estimate with exponent $1/(d+2)$, as derived in the supplement. Its dimension dependence is material.

An independent action solver must provide an \emph{upper bound} on the unobserved maximization error. Trying a few alternative actions supplies only a lower bound on the true gap. Concavity with a duality gap, interval bounds, exhaustive enumeration of a finite action set, or an action net with a known Lipschitz modulus can provide the needed upper bound. An unverified local actor optimizer cannot certify itself.

% Reviewed source: revisions/2026-09-29-r6/manuscript/bridge.tex, line 1, commit f5021cefa71492babfcbfa580e0c984f59a9de26
\subsubsection{Differential and discrete error accounts}\label{sec:bridge-r6}
For the stationary additive problem, write $f=r_0-\rho v$, with $\rho>0$, and let $R_h$ restrict a candidate to the nodes of a positive-weight scheme. On interior nodes set
\[
 T_h^a w=q_h(P_h^aw+h r_0^a),\qquad q_h=(1+\rho h)^{-1}.
\]
Stopping nodes instead use their prescribed payoff.
\begin{proposition}[Transfer to a Bellman certificate]\label{prop:bridge-r6}
Suppose the represented policy is feasible for the discrete action class and
\[
 \sup_a\left\|h^{-1}(P_h^aR_hv-R_hv)-R_h\LL^av\right\|_{\infty,\mathrm{int}}\leq\kappa_h.
\]
If the differential evaluation residual is at most $e$ and its feasible-action gap at most $\eta$, then, on the interior nodes,
\[
 \|R_hv-T_h^aR_hv\|_\infty\leq q_hh(e+\kappa_h),\qquad
 \|T_hR_hv-T_h^aR_hv\|_\infty\leq q_hh(\eta+2\kappa_h).
\]
Boundary discrepancies and finite-action coverage errors are additional accounts.
\end{proposition}
The proof follows by adding and subtracting $R_h\LL^av$; the maximization inequality uses the consistency error twice. The converse is false for arbitrary oscillatory interpolants. For smooth interior functions, local cubature and interpolation estimates can bound $\kappa_h$; they do not justify assigning zero error to a stopping or reflected boundary without checking its realization.

For finite-horizon maps with Lipschitz factors $\beta_t$, let $\delta_t$ bound evaluation error, $\epsilon_t$ bound improvement error, and $B_N$ bound terminal error. Define
\begin{equation}\label{eq:accounts-r6}
 E_N=B_N,\quad O_N=B_N,\quad R_N=2B_N,\qquad
 \begin{cases}
 E_t=\delta_t+\beta_t E_{t+1},\\
 O_t=\delta_t+\epsilon_t+\beta_t O_{t+1},\\
 R_t=2\delta_t+\epsilon_t+\beta_t R_{t+1}.
 \end{cases}
\end{equation}
Backward comparison bounds policy-value error by $E_t$, optimal-value error by $O_t$, and policy regret by $R_t$. These arrays are computed, rather than replaced by a training loss, in the NDU experiment. Bounds at later time indices must use their own accounts, not necessarily the account at time zero. Representation and incomplete optimization enter $\delta_t$; action coverage enters $\epsilon_t$; the error of replacing the controlled diffusion by a scheme is separate. If a uniform continuous-to-discrete policy-payoff discrepancy is bounded by $D_h$, transferring a discrete regret bound adds $2D_h$. Without that estimate, the certificate remains a certificate for the stated discrete economy. In particular, this accounting does not establish the optimizer rate required by Theorem~\ref{thm:viscosity}.
% Reviewed source: revisions/2026-10-03-r8/manuscript/method.tex, line 1, commit f5021cefa71492babfcbfa580e0c984f59a9de26
\subsubsection{Continuous-action proposals and a budget-safe Bellman enclosure}\label{sec:r8-operator}
The learned map and its verifier have different computational roles. In the differential implementation, a feasible actor supplies an action and the critic evaluates that action's generator. In the positive-weight implementation, write
\begin{equation}\label{eq:r8-positive}
 (T_n^a w)_i=r_{n,i}(a)+q_n\{P_{n,i}(a)w+b_{n,i}(a)\},
 \qquad (T_nw)_i=\sup_{a\in\Act_i}(T_n^aw)_i .
\end{equation}
Here $P_{n,i}(a)$ has nonnegative entries with sum at most one. The term $b$ contains prescribed stopping payoffs, and $q_n>0$. Continuous controls need not enter this expression concavely. Reflection, interpolation cells, and stopping decisions can change as the action changes. A supremum, rather than an attained maximizer, suffices below.

Our continuous-action implementation fits the actor by differentiating its own evaluated continuation payoff. It does not construct an all-action target or receive an exact maximizing label. A bounded tanh output enforces the control box. A local improvement step can subsequently compare that proposal with finitely many continuous perturbations. Such a step is recorded as part of a \emph{hybrid} policy; it is not identified with the raw network and is not assumed globally optimal. Critic fitting uses the selected feasible action and the preceding continuation critic. Training ends before any exhaustive reference or action-cover verification begins.

For independent verification, freeze the complete policy $\pi=(\pi_n)$ and cover each action box by closed subboxes. At a state and time, let $c_i$ be an outward lower bound obtained by evaluating any feasible action using a fixed continuation upper array. A box can be discarded only when its verified upper bound does not exceed $c_i+\tau$. Every box not so discarded is either subdivided or retained with its upper bound. Exhausting a time, depth, or box budget therefore enlarges the reported bound; it never converts an unresolved search into a success.

\begin{proposition}[Budget-safe action and policy envelopes]\label{prop:r8-envelope}
Suppose the terminal arrays satisfy $L_N\leq g\leq U_N$. At each time $n$, evaluate the fixed-policy map outward to obtain $L_n\leq T_n^{\pi_n}L_{n+1}$. For every state, maintain a complete action cover whose box bounds enclose $T_n^aU_{n+1}$ for all actions in that box. Set $U_{n,i}$ to the maximum of the outward-rounded pruning threshold and all unresolved box upper bounds. Then, regardless of which boxes remain unresolved,
\begin{equation}\label{eq:r8-envelope}
 L_n\leq V_n^\pi\leq V_n^*\leq U_n,
 \qquad 0\leq V_n^*-V_n^\pi\leq U_n-L_n .
\end{equation}
All inequalities are on the stated nodal economy. No concavity, exact training labels, or convergence of neural optimization is required. If $U_n-T_nU_{n+1}\leq e_n\boldsymbol 1$, then its optimal-value excess is at most
\[
 \left(\prod_{j=n}^{N-1}q_j\right)\|U_N-g\|_\infty+
 \sum_{k=n}^{N-1}\left(\prod_{j=n}^{k-1}q_j\right)e_k.
\]
\end{proposition}

The proof, including stopping faces, is in the supplement. A stored unresolved bracket supplies an upper bound for $e_n$; it is not replaced by the requested tolerance. The policy lower envelope and the Bellman upper envelope are computed separately. The latter may be shared across several frozen policies on the same economy, while each policy requires its own evaluation. This separation prevents an expensive verification oracle from becoming undisclosed actor supervision.

In the preference application, the continuous box has three coordinates and the state grid has two. Bilinear interpolation is enclosed cell by cell. On each rectangle, its extrema occur at cell corners; a second implementation uses cell-slope enclosures. If a box crosses a wealth stopping face, the verifier takes the union of the interior continuation and the prescribed exit payoff. Their numerical values need not coincide at off-grid exit locations. Treating the entire action objective as globally differentiable at this switch would invalidate a derivative-only bound. Reflection is handled by including every crossed fold. Positive consumption and a preference interval strictly above one keep the powers and denominators in their stated domains.

The complete-action bound is not a continuous-state or continuous-time theorem. The verifier interprets stored primitive coefficients and nodal abscissae as exact real inputs and encloses subsequent operations outward. Training and verification agree to numerical precision on the inherited action grid, but the certificate's precise object is this declared nodal economy, not an unspecified sequence of floating-point evaluations. The supplemental transfer result states separately the state interpolation, diffusion approximation, boundary monitoring, and action errors needed to recover the continuous economy. None is inferred from a sampled residual or silently assigned zero.
